\section{Literature Review}

\subsection{Static Android Exposure and Analysis Limitations}
Static Android research has measured packaged exposure at scale, but the cited studies do not share a unit of analysis. ManiScope performs manifest-level structural validation across Google Play and pre-installed apps, providing reproducible static evidence of misconfiguration without a corresponding runtime characterization \cite{yang2022maniscope}. Wang et al.\ instead examine runtime-permission handling, showing that realized permission behavior depends on Android version and execution state and therefore cannot be read from declarations alone \cite{wang2023runtimepermissions}. Zhu et al.\ shift the unit of analysis to SAST tools themselves: vulnerability-type coverage and reporting semantics vary across detectors, so raw finding counts are tool-dependent exposure evidence rather than universal vulnerability truth \cite{zhu2024sastandroid}. Khedkar et al.\ further show that privacy-related data collection can arise through system APIs, user interfaces, and third-party code, which complicates a purely manifest-level interpretation \cite{khedkar2026privacydata}. Taken together, static evidence is scalable, versionable, and useful for identifying packaged exposure, but its meaning depends on detector semantics and the analyzed surface, and it does not establish whether a capability becomes active during execution of the same identified build.

\subsection{Runtime and Network Observation}
Sutter et al.\ review dynamic Android analysis at the methodology level, identifying coverage, stimulation strategy, nondeterminism, and environment dependence as recurring limits on reproducibility and generalization \cite{sutter2024dynamic}. Paci et al.\ provide an empirical counterpart: execution-time network evidence can reveal third-party tracking infrastructure that static inspection alone may not expose \cite{paci2023tracking}. Sutter et al.\ therefore explain why runtime measurements are difficult to reproduce, while Paci et al.\ demonstrate why runtime network evidence is privacy-relevant. Runtime-only analysis, however, generally lacks the same-build static context needed to interpret packaged exposure and observed behavior together.

\subsection{Privacy Disclosures and Third-Party SDK Behavior}
Arkalakis et al.\ dynamically compare measured app behavior with Google Play Data Safety disclosures and report widespread incomplete disclosure \cite{arkalakis2024datasafety}. Khandelwal et al.\ study the same label system at larger scale from both measurement and developer perspectives, documenting under-reporting, over-reporting, and process friction \cite{khandelwal2024privacylabels}. Zhang et al.\ move lower in the software supply chain to privacy-configurable SDK APIs and the implementation risks of developer-facing privacy controls \cite{zhang2024privacysdks}. Rodriguez et al.\ similarly show that third-party SDK settings and defaults can shape downstream application privacy posture \cite{rodriguez2025privacysettings}. Disclosure studies ask whether declared practices correspond to observed behavior, whereas SDK studies examine how third-party configuration influences the app itself. Both provide important context, but neither replaces exact-build technical evidence connecting packaged exposure with runtime observations.

\subsection{Combined Static-Runtime Studies}
Karyotakis et al.\ are the closest combined precedent: they compare static and dynamic security and privacy characteristics of three Android messaging apps under controlled, reproducible scenarios \cite{karyotakis2026messaging}. The studies share the motivation of combining evidence layers, but they differ in scope and provenance. Karyotakis et al.\ focus deeply on a single functional class, while the present study emphasizes selected-build alignment across 15 heterogeneous consumer apps, preserves package/version/hash identity, reports strict-idle, QFG, and interactive classes separately, and uses a fixed recent-evidence window. A build may declare permissions, exported components, and SDK configuration that a short runtime scenario does not fully exercise, while runtime captures can reveal service activation that is not visible from manifest inspection. The comparison is bounded: this study is not an ecosystem-prevalence measurement and remains single-device and scenario-bounded. Table~\ref{tab:literature-comparison} summarizes that positioning.

\begin{table*}[t]
\centering
\caption{Comparison of related Android privacy and security studies.}
\label{tab:literature-comparison}
\tiny
\begin{tabular}{@{}p{0.18\textwidth}p{0.075\textwidth}p{0.075\textwidth}p{0.10\textwidth}p{0.125\textwidth}p{0.115\textwidth}p{0.15\textwidth}@{}}
\toprule
Study & Static & Runtime & Privacy & Provenance & Scope & Limit \\
\midrule
ManiScope \cite{yang2022maniscope} & Yes & No & Limited & Manifest/package & Large app corpus & No runtime alignment \\
Runtime permissions \cite{wang2023runtimepermissions} & Permissions & Behavioral failures & Indirect & Version/context varies & Permission issues & Not network-focused \\
Dynamic SLR \cite{sutter2024dynamic} & Surveyed & Surveyed & Varies & Method-level & Literature & No empirical app bundle \\
Tracking study \cite{paci2023tracking} & Some & Yes & Tracking & Execution-based & Third-party tracking & Less build alignment \\
Data Safety longitudinal \cite{arkalakis2024datasafety} & Ecosystem & Dynamic & Disclosure & Longitudinal & Google Play labels & Disclosure-centered \\
Privacy labels \cite{khandelwal2024privacylabels} & Measurement & No & Disclosure & Developer/app & Data Safety & No runtime PCAP \\
Privacy SDKs \cite{zhang2024privacysdks} & SDK config & Limited & SDK risk & SDK/app & Mobile SDKs & SDK-specific \\
Android SAST \cite{zhu2024sastandroid} & Yes & No & Security findings & Tool/app & SAST tools & No runtime layer \\
This study & Yes & Yes & Transport/provenance & Exact selected build & 15 apps & Scenario-bounded \\
\bottomrule
\end{tabular}
\end{table*}


\subsection{Synthesis and Research Gap}
Static research provides breadth and structural evidence, but not realized same-build behavior. Runtime and tracking studies reveal execution-time communication, yet they are scenario-dependent and often lack build-locked static context. Disclosure and SDK work operate at policy and configuration layers and answer related but different questions. Combined studies narrow the gap, but existing work is often smaller in functional scope or does not preserve exact-build provenance and QFG separation. The present study therefore treats the selected app build as the empirical unit: package, version, and APK-hash alignment; static and runtime evidence from that same build; separate strict-idle, QFG, and interactive classes; and disagreement between layers retained as risk-posture information rather than collapsed into a score. Those controls improve reconstructability of the reported rows without converting a single 14-day window into a longitudinal experiment. This remains a single-device, scenario-bounded study.
